\section[\nierozwiazane{Zadanie 1}]{Zadanie 1}
Rozważ dwustanowy proces \(\mathcal{X} = (X_t)_{t \in \mathbb{N}}\), który jest definiowany przez
dwie macierze \(P\) oraz \(Q\), obie rozmiaru \(2 \times 2\). W parzystych krokach macierz \(P\) definiuje
przejścia pomiędzy stanami, zaś w nieparzystych macierz \(Q\). Uzasadnij dlaczego ten proces
niekoniecznie jest łańcuchem Markowa. Znajdź analogiczny proces, ale z większą liczbą
stanów, który już jest czasowo homogeniczny, a zatem jest łańcuchem Markowa, tzn.
znajdź proces \(\mathcal{Y} = (Y_t)_{t \in \mathbb{N}}\) oraz funkcję \(f\) ze zbioru stanów procesu \(\mathcal{Y}\) w zbiór stanów \(\mathcal{X}\)
takie, że (1) \(\mathcal{Y}\) jest łańcuchem Markowa, (2) \((f(Y_t))_{t \in \mathbb{N}}\) jest takim samym procesem co \(\mathcal{X}\).

\brakrozwiazania
